\icode{E-MD-SI-L2-261113-004}
如图为粮袋的传送装置，已知 \(AB\) 间长度为 \(L\)，传送带与水平方向的夹角为 \(\theta\)，工作时其运行速度为 \(v\)，方向如图所示，粮袋与传送带间的动摩擦因数为 \(\mu\)，正常工作时工人在 \(A\) 点将粮袋放到运行中的传送带上，关于粮袋从 \(A\) 到 \(B\) 的运动，以下说法正确的是（最大静摩擦力等于滑动摩擦力，重力加速度为 \(g\)）\paren[]
	
\begin{tikzpicture}[>=stealth,font=\scriptsize]
	\begin{scope}[rotate=37]
		\coordinate (o) at(0,0);
		\coordinate (length) at(3,0);
		\coordinate (square) at(0,.3);
		\draw[gray] (0,0) circle(.3) (length) circle(.3);
		\draw[gray] (square)--++(0,.5) (3,.3)--++(0,.5);
		\draw[<->] (0,.6)node[left]{\(A\)}--node[fill=white,rotate=37]{\(L\)}(3,.6)node[right]{\(B\)};
		\draw ($(o)+(0,.3)$)--($(length)+(0,.3)$) arc(90:-90:.3)--($(length)+(0,-.3)$)--($(o)+(0,-.3)$) arc(270:90:.3);
		\fill (o) circle(.07) (length) circle(.07);
		\draw[->] (1.2,.15)--node[below,rotate=37]{\(v\)}(2,.15);
	\end{scope}
	\draw[dashed,cyan] (0,-.3)--(2,-.3);
	\draw[gray] ({.3*tan(18.5)+.4},-.3) arc(0:37:.4);
	\node at(.6,-.12){\(\theta\)};
\end{tikzpicture}
\begin{choices}
	\item 动摩擦因数 \(\mu\) 一定大于 \(\tan\theta\)
	\item 动摩擦因数 \(\mu\) 可能小于 \(\tan\theta\)
	\item 粮袋开始运动的加速度为 \(g(\mu\cos\theta + \sin\theta)\)
	\item 粮袋到达 \(B\) 点的速度可能大于 \(v\)
\end{choices}
\begin{answer}
    A\qquad 提示：
\end{answer}